% --- scholar LaTeX fragment --------------------------------------------
% Generated by: scholar sessions export -f table
% Session: ovl-property
% \input this file from the body of a document whose
% preamble loads:
%   \usepackage{longtable}
%   \usepackage{booktabs}
%   \usepackage{array}
% ----------------------------------------------------------------------

\begingroup\footnotesize
\begin{longtable}{@{}%
  >{\raggedright\arraybackslash}p{0.44\textwidth}c%
  >{\raggedright\arraybackslash}p{0.13\textwidth}%
  >{\raggedright\arraybackslash}p{0.26\textwidth}@{}}
\caption{Träffar som rör påståendet, Egenskaper och enhetlig åtkomst (6 citerade, 0 som stöder påståendet och 0 som kvalificerar eller motsäger det, av 847 unika träffar; de 502 angränsande och 339 felträffarna listas inte). Skälet till varje inkludering står i sista kolumnen; för maskinklassade rader anges modellens konfidens. Databaser: OpenAlex (OA; inte open access), IEEE Xplore (IEEE), Web of Science (WoS), Scopus.}\label{tab:sok-egenskaper}\\
\toprule
Titel & År & Leverantör & Skäl \\
\midrule
\endfirsthead
\multicolumn{4}{@{}l}{\emph{\tablename~\thetable{} (forts.)}}\\
\toprule
Titel & År & Leverantör & Skäl \\
\midrule
\endhead
\midrule \multicolumn{4}{r@{}}{\emph{forts.\ på nästa sida}}\\
\endfoot
\bottomrule
\endlastfoot
On the problem of uniform references to data structures & 1975 & OA & \emph{citerad}: principens ursprung \\
On the use of properties in Java applications & 2010 & Scopus & \emph{citerad}: läsmetoder i praktiken, mätt \\
Polymorphic identifiers & 2013 & OA & \emph{citerad}: principens gränser \\
Polymorphic identifiers: Uniform resource access in objective-smalltalk & 2014 & Scopus & \emph{citerad}: principens gränser \\
Right and wrong: ten choices in language design & 2022 & OA & \emph{citerad}: läsbarhet, påstådd av upphovsmannen \\
Truly abstract interfaces for algebraic data types: The extractor typing problem & 2018 & Scopus & \emph{citerad}: principen som designvärde \\
\end{longtable}
\endgroup
